\section{A Model Problem}
\label{sec:ginzburg_landau}

The asymptotic center-manifold derived in section~\ref{sec:center_manifold_derivation} is applied to a Ginzburg-Landau model described by a complex-conjugate pair of equations as,
\begin{subequations}
	\label{eqn:ginzburg_landau}
	\begin{align}
		\dfrac{\partial A}{\partial t} =& \left[\delta_{1}\nabla + \delta_{2} + \delta_{3}x + \delta_{4}\Lap\right]A + \delta_{5}|A|^{2}A, \\
		%
		\dfrac{\partial A^{*}}{\partial t} =& \left[\delta^{*}_{1}\nabla + \delta^{*}_{2} + \delta^{*}_{3}x + \delta^{*}_{4}\Lap\right]A^{*} + \delta^{*}_{5}|A|^{2}A^{*}. 
	\end{align}
\end{subequations}
Where, $x$ is the spatial coordinate, $\nabla$ and $\Lap$ are the one-dimensional gradient and Laplacian operators, and the superscript $^{*}$ represents complex conjugation. Typically such a model would arise as the simplified dynamics of a more detailed dynamical system. Here the Ginzburg-Landau model itself is treated as the underlying dynamical model, which represents the standard system in our context. Such models have often been utilized to investigate local and global bifurcations in spatially developing flows \citep{chomaz88,chomaz05}. It is easily seen that $A(x)=0$ represents a fixed point of the system, with the linear operators given by
\begin{align}
	\mathcal{L}_{0}	=& \left[\delta_{1}\nabla + \delta_{2} + \delta_{3}x + \delta_{4}\Lap\right], \nonumber \\
	%
	\mathcal{L}^{*}_{0} =&\left[\delta^{*}_{1}\nabla + \delta^{*}_{2} + \delta^{*}_{3}x + \delta^{*}_{4}\Lap\right]. \nonumber
\end{align}

For semi-infinite domains, $x\in[0,\infty)$, the spectrum for $\mathcal{L}_{0}$ is known analytically \citep{chomaz88} and is given by,
%\begin{subequations}
	\begin{align}
		\omega_{j}	=& \delta_{2} - \dfrac{\delta_{1}^{2}}{4\delta_{4}} + (\delta_{4}\delta_{3}^{2})^{1/3}\zeta_{j}, \label{eqn:eigenvalues} 
		%
		%\phi_{j}(x)		=& \mathrm{Ai}([-\delta_{3}/\delta_{4}]^{1/3}x + \zeta_{j}), \label{eqn:eigenvectors}
	\end{align}
%\end{subequations}
and the corrresponding complex conjugate for $\mathcal{L}^{*}_{0}$. Here $\zeta_{j}$ denotes the countable set of zeros of the Airy function. The system parameters are chosen \emph{ad hoc} and are given in table~\ref{table:GL_bifurcation_params}, which results in a single critical eigenmode with eigenvalue $\omega_{1} = 1.0\iota$ for $\mathcal{L}_{0}$ and its conjugate, $-1.0\iota$ for $\mathcal{L}^{*}_{0}$. All other eigenvalues have negative real parts. 
\begin{table}[h]
	\renewcommand{\arraystretch}{1.5}
	\begin{tabular}{cc|cc}
			Parameter & Value & Parameter & Value \\
			\hline
			\hline
			$\delta^{b}_{1}$ & $-0.1$  & $\delta^{b}_{4}$ & $(1.0 - \iota)/\sqrt{2}$ \\
			\hline
			%
			 $\delta^{b}_{2}$ & $0.741 + 1.025\iota$ & $\delta^{b}_{5}$ & $-0.1+0.1\iota$ \\
			 \hline
			%
			$\delta^{b}_{3}$ & $-0.125$ &  $\omega_{1}$ & $+1.0\iota$ \\
			\hline	
			\hline
		\end{tabular}
		\caption{Parameters for the Ginzburg-Landau system at bifurcation.}
		\label{table:GL_bifurcation_params} 
\end{table}
 The linear operator at the bifurcation point for the state vector, $\vecu = [A; A^{*}]$ (with $\vecu\in\mathds{C}^{N}$), is written as
\begin{eqnarray}
	\Lu = \begin{bmatrix}
		\mathcal{L}_{0}		& 0 \\
		0								  & \mathcal{L}^{*}_{0}
	\end{bmatrix},
\end{eqnarray}
which, in our context, represents the linear operator for the standard system. 
 
 A center-manifold approximation of the model problem requires the evaluation of restricted resolvents, which is a difficult problem to solve analytically. Therefore the problem is considered in a trucated domain of $x\in [0,L]$ with the boundary conditions of $A = 0$ at $x=0$ and, $\partial A/\partial x = 0$ (natural boundary conditions) at $x=L$. Here the domain length is specified as $L=40$. Due to the spatially localized nature of the first few eigenvectors (which dictate the global dynamics), this truncation does not affect the spectrum in any significant manner. This can be verified from the comparison of the numerically evaluated spectrum (via the implicitly restarted Arnoldi method \citep{sorensen92}) in the truncated domain to the analytical spectrum in the semi-infinite domain (figure~\ref{fig:GL_spectrum}). 
\begin{figure}
	\centering
	\includegraphics[width=0.49\textwidth]{spectra}
	\caption{Comparison of the numerically evaluated spectrum of the Ginzburg-Landau model $(\mathcal{L}_{0})$ in the truncated domain with the analytical spectrum in the semi-infinite domain.}
	\label{fig:GL_spectrum}
\end{figure}

The standard system then has $n=2$ critical eigenvalues and, the matrices for the critical eigenvalue and eigenvectors are given as,
\begin{align}
	\LambdaC = \begin{bmatrix}
		\iota & 0 \\
		0					  & -\iota
	\end{bmatrix}; &&
	%
	\DevMatrix = [\Devec_{1}; \Devec_{2}]; &&
	%
	 \AevMatrix = [\Aevec_{1}; \Aevec_{2}], \nonumber
\end{align}
where the eigenvectors have the following form in terms of the two components,
\begin{align}
	\Devec_{1} = [\phi_{1}; \veczero]; && \Devec_{2} = [\veczero; \phi^{*}_{1}], \nonumber \\
	\Aevec_{1} = [\chi_{1}; \veczero]; && \Aevec_{2} = [\veczero; \chi^{*}_{1}], \nonumber
\end{align}
The numerically calculated eigenvectors $\phi_{1}$ and $\chi_{1}$ are shown in figure~\ref{fig:extended_eigenvectors}. 
\begin{figure}
	\centering
	\includegraphics[width=0.49\textwidth]{extended_eigenvectors}
	\caption{The numerically evaluated eigenvector $\phi_{1}$, the adjoint eigenvector $\chi_{1}$ and, the shape of the harmonic forcing $\xi$.}
	\label{fig:extended_eigenvectors}
\end{figure}

The system in equation~\eqref{eqn:ginzburg_landau} has a total of five parameters (ten including the complex-conjugate pairs) that can be varied from their bifurcation values. In order to avoid an explosion in the number of terms $T^{k}_{m}$ arising from the combinatorial problem, only the variations in $\delta_{4}$ and $\delta^{*}_{4}$ are considered. This suffices to illustrate the construction of the asymptotic center-manifold under parameter variations. Although the parameters are complex-conjugate pairs, it turns out to be simplest to treat them as independent parameters and enforce the complex-conjugate property later when parameter values are prescribed, which is essentially the same as treating $A$ and $A^{*}$ independently and having two independent critical eigenvectors $\vecv_{1}$ and $\vecv_{2}$. 
The extended variable for the parameter evolution is then $\vecnu = [\delta'_{4};\delta'^{*}_{4}]$, with 
\begin{align}
	\delta_{4} =& \delta^{b}_{4} + \delta'_{4}, \nonumber & 
	\delta^{*}_{4} =& (\delta^{b}_{4})^{*} + \delta'^{*}_{4}. \nonumber
\end{align}

Harmonic forcings will be considered to illustrate the two scenarios - resonant and non-resonant external forcings. In all cases, the shape of the external forcing is specified based on,
\begin{subequations}
	\label{eqn:forcing_shape0}
\begin{align}
	\xi(x) = \dfrac{1}{\beta}e^{-(x - 5.0)^{2}}, \\
	\beta = \left[\int_{0}^{L} |e^{-(x - 5.0)^{2}}|^{2}dx\right]^{1/2}.
\end{align} 
\end{subequations}
The forcing shape is also shown in figure~\ref{fig:extended_eigenvectors}. In principle all cases of parameter perturbations and harmonic forcings can be considered together however, in order to keep the discussion and the number of combinatorial terms manageable, they are treated as individual cases. 

\subsection{Case I: Parameter Perturbations}
\label{sec:parameter-perturbation}

First we consider the case of a pair of complex conjugate parameter perturbations $\delta'_{4}, \delta^{'*}_{4}$ without any external forcings. This results in a nonlinear system given by,
\begin{subequations}
	%\label{eqn:ginzburg_landau_param_pert}
	\begin{align}
		\dfrac{\partial A}{\partial t} =& \mathcal{L}_{0}A + \delta'_{4}\nabla^{2}A + \delta^{b}_{5} (A^{*}A)A, \nonumber \\
		%
		\dfrac{\partial A^{*}}{\partial t} =& \mathcal{L}^{*}_{0} A^{*} + \delta^{'*}_{4}\nabla^{2}A^{*} + (\delta^{b}_{5})^{*}(A^{*}A)A^{*}. \nonumber
	\end{align}
\end{subequations}
Treating $\delta^{'}_{4}$ and $\delta^{'*}_{4}$ as dynamic variables results in an extended system with $n=2$, $p=2$, $h=0$, and therefore a total critical subspace dimension of $m=4$. The eigenvalue matrices are given by,
\begin{align}
	\LambdaC = 
	\begin{bmatrix}
		\iota & 0 \\ 
		0 &  -\iota 
	\end{bmatrix}, &&
	% 
	\LambdaP = 
	\begin{bmatrix}
		0 & 0 \\ 
		0 &  0 
	\end{bmatrix}. \nonumber
\end{align}
The parameter eigenvalues are non-resonant and therefore the coupling matrix $\GammaP$ is trivial, \ie,
\begin{align}
	\GammaP = 
	\begin{bmatrix}
		0 & 0 \\ 
		0 &  0 
	\end{bmatrix}. \nonumber 
\end{align}
In this particular case since the equilibrium state at the bifurcation point is trivial $(A=\veczero)$, the linear matrix $\Lnu$ which couples $A$ and $A^{*}$ to the extended variables is also a null matrix. Therefore, through the application of algorithm~\ref{algo:extended_direct_tangent}, the critical subspace is trivially extended as,
\begin{align}
	(\lambda^{P}_{1} - \Lu)\Devec^{P}_{1} &= \veczero, \nonumber \\
	%
	(\lambda^{P}_{2} - \Lu)\Devec^{P}_{2} &= \veczero, \nonumber \\
	%
	\DevMatrix_{P} &= [\Devec^{P}_{1}, \Devec^{P}_{2}], \nonumber 
\end{align}
resulting in the final extended matrices as,
\begin{align}
	%
		\Khat =
		\begin{bmatrix}
			\LambdaC & \GammaP \\ 
			\mathbf{0} &  \LambdaP 
		\end{bmatrix}, &&
		%
		\Vhat =
		\begin{bmatrix}
			\DevMatrix & \DevMatrix_{P} \\ 
			\mathbf{0} &  \mathbf{I}_{p}
		\end{bmatrix}, \nonumber 
\end{align}
where, both $\GammaP$ and $\DevMatrix_{P}$ are null matrices (of appropriate sizes). 

Once $\EDevMatrix$ and $\Khat$ are known, algorithm~\ref{algo:normal_form} may now be applied to the extended system to obtain the center-manifold approximation to the desired order. Since $\GammaP$ is a trivial matrix $\Khat$ is diagonal and, all asymptotic vectors $\AsCenM_{i_{1},\ldots,i_{k}}$ can be evaluated independent of each other. For $m=4$, one obtains $T^{2}_{4} = 10$ terms at second order and $T^{3}_{4}=20$ terms at third order. The final reduced equations (reduced representation $\Param(\vecparv)$) are found to be,
\begin{align}
	\begin{aligned}
		\dfrac{d \parv_{1}}{d t} =&
		+\iota\parv_{1} - (0.169 - 0.153\iota)\parv_{1}\parv_{3} \\
		& -(0.008-0.014\iota)\parv_{1}\parv_{1}\parv_{2} \\
		& + (0.194 - 0.114\iota)\parv_{1}\parv_{3}\parv_{3}, \\
		%
		\dfrac{d \parv_{3}}{d t} =& 0,
	\end{aligned} \nonumber
\end{align}
along with the complex conjugated equations for $\parv_{2}$ and $\parv_{4}$. One may substitute $\parv_{1} \mapsto \parv$, $\parv_{2}\mapsto\parv^{*}$, $\parv_{3}\mapsto\delta'_{4}$, $\parv_{4}\mapsto\delta'^{*}_{4}$, to obtain expressions directly relating mode amplitude evolution to the parameter perturbations as,
\begin{align}
	\label{eqn:parameter_normal_form}
	\begin{aligned}
		\dfrac{d \parv}{d t} =&
		+\iota\parv - (0.169 - 0.153\iota)\delta'_{4}\parv \\
		& -(0.008-0.014\iota)|\parv|^{2}\parv \\
		& + (0.194 - 0.114\iota)(\delta'_{4})^{2}\parv. 
	\end{aligned}
\end{align}

The results of the asymptotic center-manifold approximations compared with the fully non-linear response of the Ginzburg-Landau system are shown in figure~\ref{fig:parameter_response}. The results compare very well for small $\delta'_{4}$ values, with progressive loss of accuracy for larger parameter variations which is typical of asymptotic results. 
\begin{figure}
	\centering
	\includegraphics[width=0.49\textwidth]{diffusion_response}
	\caption{Spatial shape of the system response $(|A|)$ to perturbations of $\delta_{4}$ parameter. The solid lines indicate results of the non-linear simulations of the Ginzburg-Landau system and the dashed lines with markers indicate the results of the center-manifold approximation.}
	\label{fig:parameter_response}
\end{figure}

A similar comparison for the angular frequency of the saturated limit-cycle is shown in figure~\ref{fig:parameter_omega}. The center-manifold approximation provides a fairly accurate reproduction of the non-linear angular frequency variation.
\begin{figure}
	\centering
	\includegraphics[width=0.49\textwidth]{diffusion_omega}
	\caption{Variation of the final limit-cycle-oscillation angular frequency due to perturbations of $\delta_{4}$ parameter. The empty circles indicate results of the non-linear simulations of the Ginzburg-Landau system and the filled circles indicate the results of the center-manifold approximation.}
	\label{fig:parameter_omega}
\end{figure}

\subsection{Case II: Non-Resonant Forcing}
\label{sec:non-resonant}

For the second case the system is forced with a pair of non-resonant complex-conjugate forcings ($h = 2$) with $\boldsymbol{\lambda}^{H} = [1.3\iota; -1.3\iota]$ and forcing shapes $\fieldf^{H}_{i}$ given by
\begin{align}
	\label{eqn:forcing_shapes}
	\fieldf^{H}_{1} = \begin{bmatrix}
			\xi(x) \\
			\veczero
		\end{bmatrix}, && \fieldf^{H}_{2} = \begin{bmatrix}
		\veczero \\
		\xi^{*}(x)
		\end{bmatrix}, 
\end{align}
where, $\xi(x)$ is defined in equation~\eqref{eqn:forcing_shape0}. The governing equations can be written as
\begin{subequations}
	\label{eqn:ginzburg_landau_forced}
	\begin{align}
		\dfrac{\partial A}{\partial t} =& \mathcal{L}_{0}A + \delta^{b}_{5} (A^{*}A)A + \xi(x)\theta,  \\
		%
		\dfrac{\partial A^{*}}{\partial t} =& \mathcal{L}^{*}_{0} A^{*} + (\delta^{b}_{5})^{*}(A^{*}A)A^{*} + \xi^{*}(x)\theta^{*},  \\
		%
		\dfrac{d\theta}{d t} =& \lambda^{H}_{1}\theta,  \\ 
		%
		\dfrac{d\theta^{*}}{d t} =& (\lambda^{H}_{1}\theta)^{*}. 
	\end{align}
\end{subequations}

No parameter perturbations are considered, resulting in $n=2$, $p=0$, $h=2$ and $m=4$. Since the external forcing is non-resonant, one obtains the following matrices,
\begin{align}
	\LambdaC = 
	\begin{bmatrix}
		\iota & 0 \\ 
		0 &  -\iota 
	\end{bmatrix}, &
	% 
	\LambdaH = 
	\begin{bmatrix}
		1.3\iota & 0 \\ 
		0 &  -1.3\iota 
	\end{bmatrix}, &
	%
	\GammaH = 
	\begin{bmatrix}
		0 & 0 \\ 
		0 &  0 
	\end{bmatrix}. \nonumber 
\end{align}
The forcing modes are non-trivial and are evaluated as follows,
\begin{align}
	(\lambda^{H}_{1} - \Lu)\Devec^{H}_{1} &= \fieldf^{H}_{1}, \nonumber \\
	%
	(\lambda^{H}_{2} - \Lu)\Devec^{H}_{2} &= \fieldf^{H}_{2}, \nonumber \\
	\DevMatrix_{H} &= [\Devec^{H}_{1}, \Devec^{H}_{2}]. \nonumber
\end{align}
The resulting extended matrices are given by
\begin{align}
	%
	\Khat =
	\begin{bmatrix}
		\LambdaC & \GammaH \\ 
		\mathbf{0} &  \LambdaH
	\end{bmatrix}, &&
	%
	\Vhat =
	\begin{bmatrix}
		\DevMatrix & \DevMatrix_{H} \\ 
		\mathbf{0} &  \mathbf{I}_{h}
	\end{bmatrix}, \nonumber 
\end{align}
where $\GammaH$ is a trivial matrix and therefore $\Khat$ is diagonal.

As done in section~\ref{sec:parameter-perturbation}, the application of algorithm~\ref{algo:normal_form} results in the evolution equations in their normal form for the reduced system,  
\begin{align}
	\label{eqn:nonresonant_normal_form}
	\begin{aligned}
		\dfrac{d \parv}{d t} =&
		+\iota\parv \\
		& -(0.008-0.014\iota)|\parv|^{2}\parv \\
		& - (0.026 - 0.051\iota)|\theta|^{2}\parv, \\
		%
		\dfrac{d \theta}{d t} =& +1.3\iota\theta.
	\end{aligned}
\end{align}
Here, the substitutions $\parv_{1} \mapsto \parv$, $\parv_{2}\mapsto\parv^{*}$, $\parv_{3}\mapsto\theta$ and $\parv_{4}\mapsto\theta^{*}$, have already been performed. One can observe that the amplitude evolution of the original critical mode, denoted by $z$, only depends on the forcing amplitude $|\theta|$ and not on its instantaneous phase. The comarison of the response of the reduced system to that of the fully non-linear system is shown in figure~\ref{fig:nonresonant_response}. The response curves are nearly identical. 
\begin{figure}
	\centering
	\includegraphics[width=0.49\textwidth]{nonresonant_response}
	\caption{Spatial shape of the system response $(|A|)$ to harmonic forcing with amplitude $|\theta|$. The solid lines indicate results of the non-linear simulations of the Ginzburg-Landau system and the dashed lines with markers indicate the results of the center-manifold approximation.}
	\label{fig:nonresonant_response}
\end{figure}

Since $|\theta|$ is an apriori known quantity which remains invariant in time, it can simply be treated as another constant paramter and, the effect of non-resonant harmonic forcing on the critical mode of the system can easily be understood by looking at the linear part of the reduced system, \ie $[-0.0258|\theta|^{2} + (1 + 0.0507|\theta|^{2})\iota]\parv$. The harmonic forcing effectively increases the angular frequency, but also induces a negative growth rate, and therefore a stabilizing effect on the original critical mode. The net result is that the original critical mode eventually decays to zero and vanishes, even if it may be initially excited, and the final response is simply the forced response to the harmonic forcing. Of course this may not always be the case and the forcing can induce a destabilizing behavior of the original critical mode. 


\subsection{Case III: Resonant Forcing}
\label{sec:resonant}

As in the previous subsection we again consider the case of the system forced with a pair of complex conjugate harmonic forcings. This time however the angular frequencies of the external forcings are resonant with the internal eigenmodes of the system. Accordingly, the eigenvalues of the forcing modes is given by $\boldsymbol{\lambda}^{H} = [\iota; -\iota]$ and the eigenvalue matrices are given by,
\begin{align}
	\LambdaC = 
	\begin{bmatrix}
		\iota & 0 \\ 
		0 &  -\iota 
	\end{bmatrix}, &&
	% 
	\LambdaH = 
	\begin{bmatrix}
		\iota & 0 \\ 
		0 &  -\iota 
	\end{bmatrix} \nonumber
\end{align}
The forcing shapes are again given by equation~\eqref{eqn:forcing_shapes} and the resulting forced system is again described by equation~\eqref{eqn:ginzburg_landau_forced}. Due to resonance, the matrix $\GammaH$ is no longer trivial and its individual entries are evaluated according to algorithm~\ref{algo:extended_direct_tangent} as,
\begin{align}
	 (\vfield{\gamma}^{H}_{1})_{1} =& \Aevec_{1}^{\dagger}\fieldf^{H}_{1}, &&
	%
	 (\vfield{\gamma}^{H}_{2})_{1} = 0, \nonumber \\
	 %
	 (\vfield{\gamma}^{H}_{1})_{2} =& 0, &&
	 %
	 (\vfield{\gamma}^{H}_{2})_{2} = \Aevec_{2}^{\dagger}\fieldf^{H}_{2}, \nonumber 
\end{align}
resulting in,
\begin{align}
	\GammaH = [\vfield{\gamma}^{H}_{1}, \vfield{\gamma}^{H}_{2}] = 
	\begin{bmatrix}
		0.151-0.481\iota & 0 \\ 
		0 &  0.151+0.481\iota
	\end{bmatrix}. \nonumber 
\end{align}
The forcing modes for the extended tangent space are evaluated as,
\begin{align}
	(\lambda^{H}_{1} - \Lu)\Devec^{H}_{1} &= \fieldf^{H}_{1} - \DevMatrix\vfield{\gamma}^{H}_{1},  \nonumber \\
	%
	(\lambda^{H}_{2} - \Lu)\Devec^{H}_{2} &= \fieldf^{H}_{2} - \DevMatrix\vfield{\gamma}^{H}_{2}, \nonumber \\
	%
	\DevMatrix_{H} &= [\Devec^{H}_{1}, \Devec^{H}_{2}]. \nonumber
	%	%
\end{align}
The extended matrices being,
\begin{align}
	%
	\Khat =
	\begin{bmatrix}
		\LambdaC & \GammaH \\ 
		\mathbf{0} &  \LambdaH
	\end{bmatrix}, &&
	%
	\Vhat =
	\begin{bmatrix}
		\DevMatrix & \DevMatrix_{H} \\ 
		\mathbf{0} &  \mathbf{I}_{h}
	\end{bmatrix}, \nonumber 
\end{align}
where, $\GammaH$ is no longer trivial and $\Khat$ has an upper triangular structure.

The asymptotic center-manifold evaluation through the application of algorithm~\ref{algo:normal_form} results in the following reduced system,
\begin{align}
	\label{eqn:resonant_normal_form}
	\begin{aligned}
		\dfrac{d \parv}{d t} =&
		+\iota\parv + (0.151 - 0.482\iota)\theta \\
		& - (0.008 - 0.014\iota)|\parv|^{2}\parv \\
		& + (0.012 + 0.019\iota)\parv^{2}\theta^{*} \\
		& - (0.012 - 0.020\iota)|\parv|^{2}\theta \\
		& + (0.001 + 0.070\iota)|\theta|^{2}\parv \\
		& - (0.013 - 0.031\iota)\parv^{*}\theta^{2} \\
		& - (0.004 - 0.067\iota)|\theta|^{2}\theta \\
		%
		\dfrac{d \theta}{d t} =& + \iota\theta,
	\end{aligned}
\end{align}
and its conjugate for $\parv^{*}$ and $\theta^{*}$. Clearly, in the case of resonant forcings, the original critical mode evolution depends on the amplitude as well as the phase of the external forcing. 

A comparison of the system response between the non-linear results and the center-manifold approximation is shown in figure~\ref{fig:resonant_response}. Again a good agreement is found for small values of $|\theta|$ with systematic deviation for larger amplitudes typical of asymptotic methods. 
\begin{figure}
	\centering
	\includegraphics[width=0.49\textwidth]{resonant_response}
	\caption{Spatial shape of the system response $(|A|)$ to resonant harmonic forcing with amplitude $|\theta|$. The solid lines indicate results of the non-linear simulations of the Ginzburg-Landau system and the dashed lines with markers indicate the results of the center-manifold approximation.}
	\label{fig:resonant_response}
\end{figure}

As a final note, one may observe that performing the substitutions, $\delta'_{4} \mapsto 0$ in equation~\eqref{eqn:parameter_normal_form} and, $\theta \mapsto 0$ in equations~\eqref{eqn:nonresonant_normal_form} and \eqref{eqn:resonant_normal_form}, reduces all three equations to,
\begin{align}
	\label{eqn:standard_normal_form}
	\begin{aligned}
		\dfrac{d \parv}{d t} =&
		+\iota\parv - (0.008-0.014\iota)|\parv|^{2}\parv.
	\end{aligned}
\end{align}
This is the normal form of the center-manifold for the standard Ginzburg-Landau system \eqref{eqn:ginzburg_landau}. This provides further numerical validation of the fact that the center-manifold approximation of the extended system reduces to that of the standard system whenever the system perturbation and harmonic forcings vanish.





